\documentclass[11pt]{amsart}

\usepackage{amssymb,amsmath,amsthm}
\usepackage{graphicx}
\usepackage{color}
\usepackage{thmtools}
\usepackage{thm-restate}
\usepackage[shortlabels]{enumitem}
\usepackage{mathrsfs}
\usepackage[utf8]{inputenc}
\usepackage[T1]{fontenc}
\usepackage{dsfont}
\usepackage{soul}
\usepackage[left=1in,right=1in,top=1in,bottom=1in]{geometry}

\setlength{\textwidth}{\paperwidth}
\addtolength{\textwidth}{-2in}
\calclayout


\newtheorem{theorem}{Theorem}[section]
\newtheorem*{theorem*}{Theorem}
\newtheorem*{claim*}{Claim}
\newtheorem{lemma}[theorem]{Lemma}
\newtheorem{proposition}[theorem]{Proposition}
\newtheorem{problem}[theorem]{Problem}
\newtheorem{question}[theorem]{Question}
\newtheorem{corollary}[theorem]{Corollary}
\newtheorem{fact}[theorem]{Fact}
\newtheorem{conjecture}[theorem]{Conjecture}
\theoremstyle{definition}
\newtheorem{definition}[theorem]{Definition}
\newtheorem{example}[theorem]{Example}
\newtheorem{xca}[theorem]{Exercise}
\theoremstyle{remark}
\newtheorem{remark}[theorem]{Remark}
\newtheorem{observation}[theorem]{Observation}
\numberwithin{equation}{section}
%%%%%%%%%%%%

%%%%%%%% NUMBERS n CONSTANTS

\newcommand{\fraka}{\mathfrak{a}}
\newcommand{\frakA}{\mathfrak{A}}
\newcommand{\frakb}{\mathfrak{b}}
\newcommand{\frakc}{\mathfrak{c}}
\newcommand{\frakd}{\mathfrak{d}}
\newcommand{\frakg}{\mathfrak{g}}
\newcommand{\fraki}{\mathfrak{i}}
\newcommand{\frakk}{\mathfrak{k}}
\newcommand{\frakn}{\mathfrak{n}}
\newcommand{\frakN}{\mathfrak{N}}
\newcommand{\frakp}{\mathfrak{p}}
\newcommand{\frakr}{\mathfrak{r}}
\newcommand{\fraks}{\mathfrak{s}}
\newcommand{\frakt}{\mathfrak{t}}
\newcommand{\fraku}{\mathfrak{u}}
\newcommand{\frakx}{\mathfrak{x}}
\newcommand{\fraky}{\mathfrak{y}}
\newcommand{\frakz}{\mathfrak{z}}

\newcommand{\frakdg}{\mathfrak{dg}}
\newcommand{\frakgr}{\mathfrak{gr}}
\newcommand{\frakros}{\mathfrak{ros}}
\newcommand{\frakphil}{\mathfrak{phil}}
\newcommand{\frakschur}{\mathfrak{schur}}

\newcommand{\frakpairs}{\mathfrak{pairs}}
\newcommand{\frakspairs}{\mathfrak{spairs}}

\newcommand{\eps}{\varepsilon}

\newcommand{\zero}{\normalfont{\textbf{0}}}
\newcommand{\one}{\normalfont{\textbf{1}}}
%\newcommand{\one}{\mathbbm{1}}

\newcommand{\emptyseq}{\epsilon}

\newcommand{\omu}{\ol{\mu}}
\newcommand{\onu}{\ol{\nu}}

%%%%%%%% SETS n ALGEBRAS

\newcommand{\N}{\mathbb{N}}
\newcommand{\R}{\mathbb{R}}
\newcommand{\Q}{\mathbb{Q}}
\newcommand{\Z}{\mathbb{Z}}
\providecommand{\C}{}\renewcommand{\C}{\mathbb{C}}
\newcommand{\K}{\mathbb{K}}
\newcommand{\G}{\mathbb{G}}
\newcommand{\PP}{\mathbb{P}}

\newcommand{\aA}{\mathcal{A}}
\newcommand{\bB}{\mathcal{B}}
\newcommand{\cC}{\mathcal{C}}
\newcommand{\dD}{\mathcal{D}}
\newcommand{\eE}{\mathcal{E}}
\newcommand{\fF}{\mathcal{F}}
\newcommand{\gG}{\mathcal{G}}
\newcommand{\hH}{\mathcal{H}}
\newcommand{\iI}{\mathcal{I}}
\newcommand{\jJ}{\mathcal{J}}
\newcommand{\kK}{\mathcal{K}}
\newcommand{\lL}{\mathcal{L}}
\newcommand{\mM}{\mathcal{M}}
\newcommand{\nN}{\mathcal{N}}
\newcommand{\pP}{\mathcal{P}}
\newcommand{\qQ}{\mathcal{Q}}
\newcommand{\rR}{\mathcal{R}}
\newcommand{\sS}{\mathcal{S}}
\newcommand{\tT}{\mathcal{T}}
\newcommand{\uU}{\mathcal{U}}
\newcommand{\vV}{\mathcal{V}}
\newcommand{\wW}{\mathcal{W}}
\newcommand{\xX}{\mathcal{X}}
\newcommand{\yY}{\mathcal{Y}}
\newcommand{\zZ}{\mathcal{Z}}

\newcommand{\aAs}{\mathscr{A}}
\newcommand{\aAa}{\aAs}
\newcommand{\bBs}{\mathscr{B}}
\newcommand{\bBb}{\bBs}
\newcommand{\cCs}{\mathscr{C}}
\newcommand{\cCc}{\cCs}
\newcommand{\dDs}{\mathscr{D}}
\newcommand{\dDd}{\dDs}

%%%%%%%% COMMANDS & OPERATORS

\DeclareMathOperator{\add}{add}
\DeclareMathOperator{\ba}{ba}
\DeclareMathOperator{\card}{card}
\DeclareMathOperator{\cf}{cf}
\DeclareMathOperator{\cof}{cof}
\newcommand{\concat}{{^\smallfrown}}
\DeclareMathOperator{\cov}{cov}
\DeclareMathOperator{\der}{\mathrm{d}}
\DeclareMathOperator{\dens}{dens}
\DeclareMathOperator{\diam}{diam}
\DeclareMathOperator{\dom}{dom}
\newcommand{\embeds}{{\hookrightarrow}}
\DeclareMathOperator{\ev}{ev} \DeclareMathOperator{\Id}{Id}
\DeclareMathOperator{\Imc}{Im} \DeclareMathOperator{\intt}{int}
\newcommand{\lra}{\longrightarrow}
\newcommand{\medcup}{\scalebox{1.3}{\ensuremath{\cup}}}
\newcommand{\medvee}{\scalebox{1.3}{\ensuremath{\vee}}}
\DeclareMathOperator{\non}{non}
\newcommand{\ol}{\overline}
\let\Proj\relax\DeclareMathOperator{\Proj}{{\rm Proj}}
\DeclareMathOperator{\ram}{ram}
\DeclareMathOperator{\ran}{ran}
\DeclareMathOperator{\Rec}{Re}
\newcommand{\rstr}{\restriction}
\newcommand{\rstra}{\restriction}
\DeclareMathOperator{\sgn}{sgn}
\newcommand{\sm}{\setminus}
\newcommand{\sub}{\subseteq}
\DeclareMathOperator{\supp}{supp}
\DeclareMathOperator{\spn}{span}
\DeclareMathOperator{\clspn}{\overline{span}}%\newcommand{\clspn}{{\overline{\rm span}\,}}
\newcommand{\w}{\wedge}
\newcommand{\wh}{\widehat}

\newcommand{\seqf}[1]{\big\langle#1\big\rangle}
\newcommand{\seq}[2]{\big\langle#1\colon\ #2\big\rangle}
\newcommand{\seqn}[1]{\big\langle#1\colon\ n\io\big\rangle}
\newcommand{\seqN}[1]{\big\langle#1\colon\ N\io\big\rangle}
\newcommand{\seqk}[1]{\big\langle#1\colon\ k\io\big\rangle}
\newcommand{\seql}[1]{\big\langle#1\colon\ l\io\big\rangle}
\newcommand{\seqm}[1]{\big\langle#1\colon\ m\io\big\rangle}
\newcommand{\seqj}[1]{\big\langle#1\colon\ j\io\big\rangle}
\newcommand{\seqnk}[1]{\big\langle#1\colon\ n,k\io\big\rangle}

\newcommand{\clopen}[1]{\left[#1\right]}

\newcommand{\ctblsub}[1]{\left[#1\right]^\omega}
\newcommand{\finsub}[1]{\left[#1\right]^{<\omega}}

%%%%%%%% SPACES

\newcommand{\iA}{\in\aA}
\newcommand{\iF}{\in\fF}
\newcommand{\notiF}{\not\in\fF}

\newcommand{\io}{\in\mathbb{N}}%{\in\N}
\newcommand{\lo}{<\omega}

\newcommand{\wo}{{\wp(\omega)}}%{{\wp(\N)}}
\newcommand{\bo}{{\beta\omega}}%{{\beta\N}}
\newcommand{\oo}{\omega^\omega}%{\N^\N}
\newcommand{\Ro}{\R^\N}
\newcommand{\cso}{\ctblsub{\omega}}%{\ctblsub{\N}}
\newcommand{\fso}{\finsub{\omega}}%{\finsub{\N}}
\newcommand{\seqomega}{\omega^{<\omega}}
\newcommand{\ipo}{\in\wo}
\newcommand{\ioo}{\in\oo}

\newcommand{\elli}{{\ell_\infty}}
\newcommand{\ellid}{{\ell_\infty^*}}

\newcommand{\ooi}{(\oo)^\infty}

\newcommand{\Fro}{{Fr(\omega)}}
\newcommand{\Cantor}{2^\omega}
\newcommand{\two}{\{0,1\}}

\newcommand{\bao}{{\ba(\wo)}}

%%%%%%%% VARIA

\newcommand{\stevo}{Todor\v{c}evi\'c}

\newcommand{\noproof}{\hfill$\Box$}
\newcommand{\noproofm}{\hfill\Box}

\newcommand{\horline}{\noindent\hrulefill}

%\renewcommand{\io}{\in\omega}
\newcommand{\forces}{\Vdash}
\newcommand{\Sacks}{\mathbb{S}^\kappa}
\renewcommand{\SS}{\mathbb{S}}

\newcommand{\grp}{the Grothendieck property}
\newcommand{\gr}{Grothendieck\ }
\newcommand{\nikp}{the Nikodym property}
\newcommand{\nik}{Nikodym\ }
\newcommand{\jnp}{the Josefson--Nissenzweig property}
\newcommand{\jn}{Josefson--Nissenzweig\ }


\begin{document}

% \title[short text for running head]{full title}
\title[Complementability of $c_0$ in spaces $C(K\times L)$]{On complementability of $c_0$ in spaces $C(K\times L)$}
\author[J. K\k{a}kol]{Jerzy K\k{a}kol}
\address{Faculty of Mathematics and Computer Science, A. Mickiewicz University, Pozna\'n, Poland, and Institute of Mathematics, Czech Academy of Sciences, Prague, Czech Republic.}
\email{kakol@amu.edu.pl}
\author[D.\ Sobota]{Damian Sobota}
\address{Universit\"at Wien, Institut f\"ur Mathematik, Kurt G\"odel Research Center, Wien, Austria.}
\email{ein.damian.sobota@gmail.com}
\urladdr{www.logic.univie.ac.at/~{}dsobota}
\author[L. Zdomskyy]{Lyubomyr Zdomskyy}
\address{Universit\"at Wien, Institut f\"ur Mathematik, Kurt G\"odel Research Center, Wien, Austria.}
\email{lzdomsky@gmail.com}
\urladdr{www.logic.univie.ac.at/~{}lzdomsky}
\thanks{The research of the first  named author is supported by the GA\v{C}R project 20-22230L and RVO: 67985840. The second and third named authors were supported by the Austrian Science Fund FWF, Grants I 2374-N35, I 3709-N35, M 2500-N35, I 4570-N35.}
\subjclass[2020]
{Primary: 
46E15, %Banach spaces of continuous, differentiable or analytic functions
28A33, %Spaces of measures, convergence of measures
46B09. %Probabilistic methods in Banach space theory
Secondary: 
28C05, %Integration theory via linear functionals (Radon measures, Daniell integrals, etc., NOW()), representing set functions and measures
28C15, %Set functions and measures on topological spaces (regularity of measures, etc.)
46E27. %Spaces of measures
}

\keywords{Banach spaces of continuous functions, complementability, convergence of measures, weak* topology, Josefson--Nissenzweig theorem, Weak Law of Large Numbers}

\begin{abstract}
Using elementary probabilistic methods, in particular a variant of the Weak Law of Large Numbers related to the Bernoulli distribution, we prove that for every infinite compact spaces $K$ and $L$ the product $K\times L$ admits a sequence $\langle\mu_n\colon n\in\mathbb{N}\rangle$ of normalized signed measures with finite supports which converges to $0$ with respect to the weak* topology of the dual Banach space $C(K\times L)^*$. Our approach is completely constructive---the measures $\mu_n$ are defined by an explicit simple formula. We also show that this result generalizes the classical theorem of Cembranos and Freniche which states that for every infinite compact spaces $K$ and $L$ the Banach space $C(K\times L)$ contains a complemented copy of the space $c_0$.
\end{abstract}

\maketitle

\section{Introduction}

As usual, for a compact (Hausdorff) space $X$ we denote by $C(X)$ the Banach space of continuous real-valued functions on $X$ and by $c_0$ the Banach space of all real-valued sequences which converge to $0$, both endowed with the supremum norm.

It is an easy observation that for every infinite compact space $X$ the space $C(X)$ contains a closed linear subspace isomorphic to the space $c_0$. This subspace need not be complemented---e.g., for $X=\beta\N$, the \v{C}ech--Stone compactification of the set $\N$ of non-negative integers equipped with the discrete topology, the space $C(X)$ does not contain any complemented copies of $c_0$. However, if $X$ is a product of two infinite compact spaces, then $C(X)$ always contains a complemented copy of $c_0$, as was proved by Cembranos \cite{Cem84} and Freniche \cite{Fre84}.

\begin{theorem}[Cembranos--Freniche]\label{theorem:cemfre}
For every infinite compact spaces $K$ and $L$ the Banach space $C(K\times L)$ contains a complemented copy of the space $c_0$.
\end{theorem}

A crucial step in both of the proofs is an application of the classical Josefson--Nissenzweig theorem which asserts that for every infinite-dimensional Banach space $E$ there exists a sequence $\seqn{x_n^*}$ in the dual space $E^*$ such that $\|x_n^*\|=1$ for every $n\io$ and $x_n^*(x)\to0$ for every $x\in E$ (that is, $\seqn{x_n^*}$ is convergent to $0$ with respect to the weak* topology). Since all standard proofs of the Josefson--Nissenzweig theorem are rather intricate and non-constructive, it is hard to deduce from the proofs of Cembranos and Freniche how the constructed complemented copy of $c_0$ in a given space $C(K\times L)$ basically looks like or what properties it has. In this paper we address this issue and provide an elementary and constructive proof of the following strengthening of Theorem \ref{theorem:cemfre}.

\begin{theorem}\label{theorem:main}
For every infinite compact spaces $K$ and $L$ there is a sequence $\seqn{\mu_n}$ of normalized signed measures on $K\times L$ with finite supports which is weak* convergent to $0$.% such that $\big\|\mu_n\|=1$ for every $n\io$ and $\int_{K\times L}fd\mu_n\to0$ for every $f\in C(K\times L)$.
\end{theorem}

Note that, by the virtue of the classical Riesz representation theorem, for every infinite compact space $X$, the Josefson--Nissenzweig theorem can be expressed in the following way: there is a sequence $\seqn{\mu_n}$ of normalized signed regular Borel measures on $X$ which converges to $0$ with respect to the weak* topology of the dual space $C(X)^*$. Hence Theorem \ref{theorem:main} might be treated in the first place as a special case of the Josefson--Nissenzweig theorem. However, as said above, in contrast to the latter general result, our proof of Theorem \ref{theorem:main} is completely constructive---for every compact spaces $K$ and $L$ the measures $\mu_n$ are given by an explicit simple formula. To prove that they weak* converge to $0$ we use elementary tools from probability theory, in particular a variant of the Weak Law of Large Numbers related to the Bernoulli distribution.

The Cembranos--Freniche theorem can be deduced from Theorem \ref{theorem:main} in the following way, using some $C_p$-theory. Recall that if $X$ is a Tychonoff space, then $C_p(X)$ denotes the space of all continuous real-valued functions on $X$ equipped with the pointwise topology. Similarly, by $(c_0)_p$ we mean the space $c_0$ but endowed with the topology inherited from the product $\R^\N$. The next corollary is an immediate consequence of Theorem \ref{theorem:main} and the following result due to Banakh, K\k{a}kol, and \'{S}liwa \cite[Theorem 1]{BKS19}: a Tychonoff space $X$ admits a sequence $\seqn{\mu_n}$ of finitely supported signed measures such that $\big\|\mu_n\big\|=1$ for every $n\io$ and $\int_Xfd\mu_n\to0$ for every $f\in C_p(X)$ if and only if the space $C_p(X)$ contains a complemented copy of the space $(c_0)_p$. Of course, if $X$ is compact, then the condition that $\int_Xfd\mu_n\to0$ for every $f\in C_p(X)$ is equivalent to $\seqn{\mu_n}$ being convergent to $0$ with respect to the weak* topology of $C(X)^*$.

\begin{corollary}\label{cor:main}
For every infinite compact spaces $K$ and $L$ the space $C_p(K\times L)$ contains a complemented copy of the space $(c_0)_p$.
\end{corollary}

\noindent With an aid of the Closed Graph Theorem, the result of Cembranos and Freniche can be easily deduced from the above corollary, see Proposition \ref{prop:c0p_c0} for more details.

Theorem \ref{theorem:main} gives also applications to Grothendieck $C(X)$-spaces. Recall that a Banach space $E$ is \textit{Grothendieck} if every sequence $\seqn{x_n^*}$ in the dual space $E^*$ which is weak* convergent to $0$, converges also weakly. It was proved by Cembranos \cite[Corollary 2]{Cem84} (cf. also \cite[Proposition 5.3]{Sch82}) that, for every infinite compact space $X$, the Banach space $C(X)$ is Grothendieck if and only if $C(X)$ does not contain any complemented copy of $c_0$. The latter result and Theorem \ref{theorem:cemfre} immediately yield the following theorem of Khurana \cite[Theorem 2]{Khu78}: for every infinite compact spaces $K$ and $L$ the space $C(K\times L)$ is never Grothendieck, that is, there exists a sequence $\seqn{\mu_n}$ of normalized signed regular Borel measures on $K\times L$ which converges weak* to $0$ but not weakly. Theorem \ref{theorem:main} strengthens this observation---there must even exist such $\seqn{\mu_n}$ which consists only of finitely supported measures (note that, by the Schur property, such $\seqn{\mu_n}$ cannot be weakly convergent to $0$). This yields indeed a strengthening of Khurana's theorem as there exists a compact space $X$ such that $C(X)$ is not Grothendieck but for every separable closed subspace $Y$ of $X$ the space $C(Y)$ is Grothendieck (see e.g. \cite[Example 10]{KM20}).

Let us finish by mentioning the following problem. % kindly suggested to us by the referee.
It is well-known that the space $C(\beta\N)$ is \textit{prime}, that is, every complemented infinite-dimensional subspace of $C(\beta\N)$ is isomorphic to $C(\beta\N)$ (see \cite{Lin67}). On the other hand, $C(\beta\N\times\beta\N)$ is obviously not prime, we do not know however whether there exist complemented subspaces different than the ones listed below (cf. \cite{Can21}).

\begin{question}
Does there exist a complemented infinite-dimensional subspace of $C(\beta\N\times\beta\N)$ which is not isomorphic to any of the following spaces: $C(\beta\N\times\beta\N)$, $C(\beta\N)$, $c_0$, $c_0\oplus C(\beta\N)$, and $c_0(C(\beta\N))$?
\end{question}

\section{Proof of Theorem \ref{theorem:main}}

%Khurana \cite[Theorem 2]{Khu78}, Cembranos \cite[Corollaries 2--3]{Cem84} and Freniche \cite[Corollary 2.6]{Fre84} proved that given two infinite compact spaces $K$ and $L$ the space $C(K\times L)$ does not have the Grothendieck property. In this section we will strengthen their result by proving that the product $K\times L$ does not even have the $\ell_1$-\gr property (Theorem \ref{theorem:products_fsjnp} and Corollary \ref{cor:product_ell1_grothendieck})---our proof is totally constructive and does not require any techniques from Banach space theory. As a corollary, we obtain that the space $C_p(K\times L)$ has always the quotient isomorphic to $(c_0)_p$. 

%Schachermayer \cite[Proposition 5.3]{Sch82} and Cembranos \cite[Corollary 2]{Cem84} proved that a $C(K)$-space is Grothendieck if and only if it does not contain any complemented copy of $c_0$. This implies Cembranos' result \cite[Corollary 3]{Cem84} stating that the space $C(K\times L)$ admits a complemented copy of $c_0$ for any two infinite compact spaces $K$ and $L$. Corollary \ref{cor:product_ell1_grothendieck}, together with the characterization of Banakh, K\c akol and \'Sliwa \cite[Theorem 1]{BKS19} and Proposition \ref{prop:complement_c0p_complemented_c0}, implies the following stronger result. (THAT WILL GO THE INTRODUCTION RATHER!!!!!)

%\begin{corollary}\label{cor:product_c0}
%For every two infinite compact spaces $K$ and $L$, the space $C_p(K\times L)$ admits $(c_0)_p$ as a quotient.
%\end{corollary}

Let us start with some notation and terminology. All topological spaces considered by us are assumed to be Tychonoff. For a Tychonoff space $X$, by $Bor(X)$ we denote the Borel $\sigma$-field of $X$. A function $\mu\colon Bor(X)\to\R$ is \textit{a measure} on $X$, if $\mu$ is additive, inner regular with respect to closed sets and outer regular with respect to open sets, and has bounded variation (that is, $\|\mu\|=\sup\big\{|\mu(A)|+|\mu(B)|\colon A,B\in Bor(X), A\cap B=\emptyset\big\}<\infty$). If $\mu$ attains negative values, then we say that $\mu$ is \textit{signed}. If $\|\mu\|=1$, then we say that $\mu$ is \textit{normalized}. \textit{The support} of $\mu$ is denoted by $\supp(\mu)$. If there are finite sequences $x_1,\ldots,x_n\in X$ (with all points pairwise distinct) and $\alpha_1,\ldots,\alpha_n\in\R$ such that $\mu=\sum_{i=1}^n\alpha_i\delta_{x_i}$, where by $\delta_x$ we denote the one-point measure at a point $x\in X$, then we say that $\mu$ is \textit{finitely supported}. Note that in this case we have $\supp(\mu)=\big\{x_1,\ldots,x_n\big\}$ and $\|\mu\|=\sum_{i=1}^n\big|\alpha_i\big|$. We assume that $0\in\N$ and by $\N_+$ we denote the set of positive integers, i.e. $\N_+=\N\sm\{0\}=\{1,2,3,\ldots\}$. We identify every $n\in\N$ with the set $\{0,\ldots,n-1\}$. If $A$ is a set, then by $|A|$ we denote its cardinality.

\medskip

For every $n\io_+$
put $\Omega_n=\{-1,1\}^n$ and $\Sigma_n=n\times\{n\}$ (so
$\big|\Omega_n\big|=2^n$ and $\big|\Sigma_n\big|=n$). To simplify
the notation, we will usually write $i\in\Sigma_n$ instead of
$(i,n)\in\Sigma_n$---this should cause no confusion. Put also
$\Omega=\bigcup_{n\io_+}\Omega_n$ and
$\Sigma=\bigcup_{n\io_+}\Sigma_n$, and endow these two sets with the
discrete topology. This way, we can think of the product space
$\Omega\times\Sigma$ as a countable union of pairwise disjoint
discrete rectangles $\Omega_k\times\Sigma_m$ of size $m2^k$---the
rectangles $\Omega_n\times\Sigma_n$, lying along the diagonal, will
bear a special meaning, namely, they will be the supports of
measures from the special sequence $\seq{\mu_n}{n\io_+}$ on the space
$\beta\Omega\times\beta\Sigma$ (that is, on the product of the \v{C}ech--Stone compactifications of $\Omega$ and $\Sigma$) defined as follows ($n\io_+$):
\[\mu_n=\sum_{\substack{s\in\Omega_n\\i\in\Sigma_n}}\frac{s(i)}{n2^n}\delta_{(s,i)}.\]
Then, $\supp\big(\mu_n\big)=\Omega_n\times\Sigma_n$, so $\big|\supp\big(\mu_n\big)\big|=n2^n$,
$\big\|\mu_n\big\|=1$, and
\[\pi_i\big[\supp\big(\mu_n\big)\big]\cap\pi_i\big[\supp\big(\mu_{n'}\big)\big]=\emptyset\]
for every $n\neq n'$ and $i\in\{0,1\}$ (here $\pi_i$ denotes the projection on the $i$-th coordinate). Note that for each $n\io_+$ and any two sets $A\in\wp(\Omega)$ and $B\in\wp(\Sigma)$ we have:
\[\tag{$\dagger$}\big|\mu_n\big([A]\times[B]\big)\big|\le\frac{\big|A\cap \Omega_n\big|}{2^n}\cdot\frac{\big|B\cap \Sigma_n\big|}{n},\]
where $[A]$ and $[B]$ always denote the clopen subsets of $\beta\Omega$ and $\beta\Sigma$ corresponding in the sense of the Stone duality to $A$ and $B$, respectively---since $\beta\Omega$ and $\beta\Sigma$ are extremely disconnected, we have $[A]=\ol{A}^{\beta\Omega}$ and $[B]=\ol{B}^{\beta\Sigma}$.

\medskip

In the next proposition we will prove that the sequence $\seqn{\mu_n}$, as defined above, is weak* convergent to $0$ on the space $\beta\Omega\times\beta\Sigma$. However, before we do that, we need to provide a bit of explanation of probability tools we use in the proof. For every $n\io_+$ and $i\in n$ define the function $X_i\colon\Omega_n\to\{0,1\}$ as follows: $X_i(r)=1$ if and only if $r(i)=1$, where $r\in\Omega_n$. Put $S_n=\sum_{i=0}^{n-1}X_i$, so $S_n\colon\Omega_n\to n$ is the function computing the number of $1$'s in the argument sequence $r\in\Omega_n$. For a finite set $A\in\fso$, let $P_A$ denotes the standard product probability on $\{-1,1\}^A$ (assigning $1/2^{|A|}$ to each elementary event, i.e. $P_A(\{r\})=1/2^{|A|}$ for each $r\in\{-1,1\}^{|A|}$). Recall that for every $k\le n$ it holds:
\[P_n(S_n=k)=P_n\big(\big\{r\in\Omega_n\colon\ S_n(r)=k\big\}\big)={n\choose k}1/2^n.\]
We will need the following fact, being a variant of the Weak Law of Large Numbers, which estimates the probability that $S_n(r)$ has value ``far'' (with respect to $\eps$) from $n/2$, i.e. that ``$r$ contains \emph{significantly} more (with respect to $\eps$) $1$'s than $-1$'s, or \textit{vice versa}''.

\begin{fact}\label{fact:bollobas}
If $n\io_+$ and $\eps\in\big(0,1/12\big]$ are such numbers that $n\ge48/\eps$, then:
\[P_n\big(\big|S_n-n/2\big|\ge\eps n/2\big)\le%\frac{\sqrt{2}}{\eps\sqrt{n}}\cdot e^{-\eps^2n/6}\le
\frac{\sqrt{2}}{\eps\sqrt{n}}.\]
\end{fact}
\begin{proof}
See Bollob\'as \cite[Theorem 1.7.(i)]{Bol01}.
\end{proof}

\medskip

We are ready to prove the aforementioned auxiliary proposition.

%\begin{theorem}\label{thm:jnp_products}
%The sequence $\seqn{\mu_n}$ defined above is a JN-sequence. Consequently, $(\bo)^2$ has the JNP.
%\end{theorem}
\begin{proposition}\label{prop:fsjnp_product_omega_sigma}
The sequence $\seq{\mu_n}{n\io_+}$ defined above is convergent to $0$ with respect to the weak* topology of the dual space $C(\beta\Omega\times\beta\Sigma)^*$.% Consequently, $\beta\Omega\times\beta\Sigma$ has the fsJNP.
\end{proposition}
\begin{proof}
Since $\beta\Omega\times\beta\Sigma$ is a totally disconnected compact space, to prove that $\seq{\mu_n}{n\io_+}$ is weak* convergent to $0$ it is enough to show that it converges to $0$ on every clopen subset of the form $[A]\times[B]$, where $A\in\wp(\Omega)$ and $B\in\wp(\Sigma)$ (to see this, use e.g. the Stone--Weierstrass theorem applied to the linear subspace of simple functions defined on clopen subsets of $\beta\Omega\times\beta\Sigma$). So let us fix two such sets $A$ and $B$.% For every $n\io_+$ let $k_n=\big|B\cap\Sigma_n\big|$. % Let thus $A$ and $B$ be two fixed subsets of $\Omega$ and $\Sigma$, respectively, and, for the sake of contradiction, let us assume that $\mu_n\big([A]\times[B]\big)\ge\eta$ for some $\eta>0$ and all $n\io$ (if necessary, we go to a subsequence).

Fix $\eps\in\big(0,1/12\big]$ and put:
\[I_0=\big\{n\io_+\colon\ \big|B\cap\Sigma_n\big|<2/\eps^4\big\}\]
and
\[I_1=\N_+\sm I_0=\big\{n\io_+\colon \big|B\cap\Sigma_n\big|\ge2/\eps^4\big\}.\]
% $n\io$ such that $n\ge2/\eps^4$ (so $n\ge3/\eps$), and put:%$n\ge48/\eta$ {\color{red}($n\ge3/\eps$)}, and put:
For each $i\in\{0,1\}$ we will find $N_i\io_+$ such that for every $n\ge N_i$, $n\in I_i$, it holds
\[\big|\mu_n\big([A]\times[B]\big)\big|=\big|\mu_n\big(\big[A\cap\Omega_n\big]\times\big[B\cap\Sigma_n\big]\big)\big|<2\eps.\]

\medskip

%If for some $i\in\{0,1\}$ the set $I_i$ is finite, then let immediately $N_i=1+\max I_i$. 
We first look for $N_0$. By ($\dagger$) for every $n\in I_0$ we have:
\[\big|\mu_n\big([A]\times[B]\big)\big|\le\frac{\big|A\cap \Omega_n\big|}{2^n}\cdot\frac{\big|B\cap\Sigma_n\big|}{n}\le\frac{\big|B\cap\Sigma_n\big|}{n}<\frac{2}{n\eps^4},\]
so if $I_0$ is infinite, then there exists $N_0\io_+$ such that for every $n\ge N_0$, $n\in I_0$, we have:
\[\big|\mu_n\big([A]\times[B]\big)\big|<2\eps.\]
If, on the other hand, $I_0$ is finite, then simply set $N_0=1+\max I_0$.

\medskip

Let us now look for $N_1$. We assume that $I_1$ is non-empty---otherwise simply set $N_1=0$.  For every $n\in I_1$ define the set $\Delta_{n,\eps}$ as follows:
\[\Delta_{n,\eps}=\Big\{s\in\Omega_n\colon\ \Big|\big|\big\{i\in B\cap\Sigma_n\colon\ s(i)=1\big\}\big|-\frac{\big|B\cap\Sigma_n\big|}{2}\Big|\ge\eps\frac{\big|B\cap\Sigma_n\big|}{2}\Big\},\]
so $\Delta_{n,\eps}$ denotes the event that $s\in\Omega_n$ is ``far'' (with respect to $\eps$) from having the same numbers of $1$'s and $-1$'s when restricted to the set $B$. If we put similarly:
\[\Gamma_{n,\eps}=\Big\{s\in\{-1,1\}^{B\cap\Sigma_n}\colon\ \Big|\big|\big\{i\in B\cap\Sigma_n\colon\ s(i)=1\big\}\big|-\frac{\big|B\cap\Sigma_n\big|}{2}\Big|\ge\eps\frac{\big|B\cap\Sigma_n\big|}{2}\Big\},\]
then we trivially have:
\[\tag{$\times$}\Delta_{n,\eps}=\Gamma_{n,\eps}\times\{-1,1\}^{\Sigma_n\sm B}.\]
Using this, for every $n\in I_1$ we will estimate the following values (see (1) and (3)):
\[\tag{$*$}\big|\mu_n\big(\big[A\cap\Delta_{n,\eps}\big]\times[B\cap\Sigma_n]\big)\big|=\Big|\sum_{\substack{s\in A\cap\Delta_{n,\eps}\\i\in B\cap\Sigma_n}}\frac{s(i)}{n2^n}\Big|\]
and
\[\tag{$**$}\big|\mu_n\big(\big[A\cap\big(\Omega_n\sm\Delta_{n,\eps}\big)\big]\times[B\cap\Sigma_n]\big)\big|=\Big|\sum_{\substack{s\in A\cap(\Omega_n\sm\Delta_{n,\eps})\\i\in B\cap\Sigma_n}}\frac{s(i)}{n2^n}\Big|.\]
Note that:
\[\big|\mu_n\big([A]\times[B]\big)\big|\le\big|\mu_n\big(\big[A\cap\Delta_{n,\eps}\big]\times[B\cap\Sigma_n]\big)\big|+\big|\mu_n\big(\big[A\cap\big(\Omega_n\sm\Delta_{n,\eps}\big)\big]\times[B\cap\Sigma_n]\big)\big|,\]
so obtaining ``good'' estimations of ($*$) and ($**$) will finish the proof.

\medskip

Fix $n\in I_1$ and let us start with the estimation of ($*$). Note that $\big|B\cap\Sigma_n\big|\ge48/\eps$, so recall ($\times$) and apply Fact \ref{fact:bollobas} with the set $B\cap\Sigma_n$ instead of the set $n=\{0,\ldots,n-1\}$ to get that: %$P_n\big(\Delta_{n,\eps}\big)\le\frac{\sqrt{2}}{\eps\sqrt{n}}$ and hence:
%%%%%\[\tag{$0$}P_n\big(\Delta_{n,\eps}\big)\le\frac{\sqrt{2}}{\eps\sqrt{n}}.\]
\[\tag{$0$}P_n\big(\Delta_{n,\eps}\big)=P_{B\cap\Sigma_n}\big(\Gamma_{n,\eps}\big)\cdot P_{\Sigma_n\sm B}\Big(\{-1,1\}^{\Sigma_n\sm B}\Big)=P_{B\cap\Sigma_n}\big(\Gamma_{n,\eps}\big)\le\frac{\sqrt{2}}{\eps\sqrt{\big|B\cap\Sigma_n\big|}}.\]
%\[\tag{$0$}P_n\big(\Delta_{n,\eps}\big)\le\frac{\sqrt{2}}{\eps\sqrt{\eta n}},\]
%since $\eta\le 1$.
It holds that:
\[\Big|\sum_{\substack{s\in A\cap\Delta_{n,\eps}\\i\in B\cap\Sigma_n}}\frac{s(i)}{n2^n}\Big|\le\sum_{\substack{s\in A\cap\Delta_{n,\eps}\\i\in B\cap\Sigma_n}}\frac{1}{n2^n}=\frac{\big|A\cap\Delta_{n,\eps}\big|\cdot\big|B\cap\Sigma_n\big|}{n2^n}\le\]
\[\le\frac{\big|\Delta_{n,\eps}\big|\cdot n}{n2^n}=P_n\big(\Delta_{n,\eps}\big),\]
so by (0) we get the following estimation of ($*$):
\[\tag{$1$}\Big|\sum_{\substack{s\in A\cap\Delta_{n,\eps}\\i\in B\cap\Sigma_n}}\frac{s(i)}{n2^n}\Big|\le\frac{\sqrt{2}}{\eps\sqrt{\big|B\cap\Sigma_n\big|}}.\]

\medskip

We now estimate ($**$). For every $s\in\Omega_n\sm \Delta_{n,\eps}$ we have:
\[\big|\sum_{i\in B\cap\Sigma_n}s(i)\big|=\Big|\big|\big\{i\in B\cap\Sigma_n\colon\ s(i)=1\big\}\big|-\big|\big\{i\in B\cap\Sigma_n\colon\ s(i)=-1\big\}\big|\Big|\le\]
\[\le\Big|\big|\big\{i\in B\cap\Sigma_n\colon\ s(i)=1\big\}\big|-\frac{\big|B\cap\Sigma_n\big|}{2}\Big|+\Big|\big|\big\{i\in B\cap\Sigma_n\colon\ s(i)=-1\big\}\big|-\frac{\big|B\cap\Sigma_n\big|}{2}\Big|<\]
\[<2\cdot\eps\frac{\big|B\cap\Sigma_n\big|}{2}=\eps\big|B\cap\Sigma_n\big|,\]
so:
\[\tag{$2$}\big|\sum_{i\in B\cap\Sigma_n}s(i)\big|<\eps\big|B\cap\Sigma_n\big|.\]
Next, it holds:
\[\Big|\sum_{\substack{s\in A\cap(\Omega_n\sm\Delta_{n,\eps})\\i\in B\cap\Sigma_n}}\frac{s(i)}{n2^n}\Big|\le\frac{1}{n2^n}\Big|\sum_{\substack{s\in A\cap(\Omega_n\sm\Delta_{n,\eps})\\i\in B\cap\Sigma_n}}s(i)\Big|=\]
\[=\frac{1}{n2^n}\Big|\sum_{s\in A\cap(\Omega_n\sm\Delta_{n,\eps})}\ \sum_{i\in B\cap\Sigma_n}s(i)\Big|\le\frac{1}{n2^n}\sum_{s\in A\cap(\Omega_n\sm\Delta_{n,\eps})}\big|\sum_{i\in B\cap\Sigma_n}s(i)\big|\le\]
\[\le\frac{1}{n}\max\Big\{\big|\sum_{i\in B\cap\Sigma_n}s(i)\big|\colon\ s\in\Omega_n\sm\Delta_{n,\eps}\Big\},\]
so by (2):
\[\tag{$3$}\Big|\sum_{\substack{s\in A\cap(\Omega_n\sm\Delta_{n,\eps})\\i\in B\cap\Sigma_n}}\frac{s(i)}{n2^n}\Big|<\frac{\eps\big|B\cap\Sigma_n\big|}{n}\le\frac{\eps n}{n}=\eps.\]

\medskip

Using (1), (3), and the fact that $\big|B\cap\Sigma_n\big|\ge2/\eps^4$, we conclude that for every $n\in I_1$ we have:
\[\big|\mu_n\big([A]\times[B]\big)\big|=\Big|\sum_{\substack{s\in A\cap\Omega_n\\i\in B\cap\Sigma_n}}\frac{s(i)}{n2^n}\Big|\le\Big|\sum_{\substack{s\in A\cap\Delta_{n,\eps}\\i\in B\cap\Sigma_n}}\frac{s(i)}{n2^n}\Big|+\Big|\sum_{\substack{s\in A\cap(\Omega_n\sm\Delta_{n,\eps})\\i\in B\cap\Sigma_n}}\frac{s(i)}{n2^n}\Big|<\]
\[<\frac{\sqrt{2}}{\eps\sqrt{\big|B\cap\Sigma_n\big|}}+\eps\le\eps+\eps=2\eps.\]
It follows that if for $N_1$ we take any number from $I_1$, e.g. we set $N_1=\min I_1$, then for every $n\in I_1$, $n\ge N_1$, we have:
\[\big|\mu_n\big([A]\times[B]\big)\big|<2\eps.\]

We finish the proof by denoting $N=\max\big(N_0,N_1\big)$ and seeing that for every $n\ge N$ we obviously have the same inequality, i.e.:
\[\big|\mu_n\big([A]\times[B]\big)\big|<2\eps.\]
Since $\eps\in\big(0,1/12\big]$ is arbitrary, it holds that $\lim_{n\to\infty}\mu_n\big([A]\times[B]\big)=0$ and hence $\seq{\mu_n}{n\io_+}$ is weak* convergent to $0$.
\end{proof}

%\begin{theorem}\label{theorem:products_fsjnp}
%For every two infinite compact spaces $K$ and $L$, their product $K\times L$ has the fsJNP.
%\end{theorem}

\begin{remark}
Let us note that from the above proof we can deduce for every $n\io_+$ the following estimations of the value $\big|\mu_n\big([A]\times[B]\big)\big|$, depending only on the size of the intersection $B\cap\Sigma_n$:
\[
\big|\mu_n\big([A]\times[B]\big)\big|\le
\begin{cases}
\frac{|B\cap\Sigma_n|}{n}&,\text{ if }\big|B\cap\Sigma_n\big|<\frac{2}{\eps^4},\\
\frac{\sqrt{2}}{\eps\sqrt{|B\cap\Sigma_n|}}+\frac{\eps|B\cap\Sigma_n|}{n}&,\text{ if }|B\cap\Sigma_n|\ge\frac{2}{\eps^4}.
\end{cases}
\]
\end{remark}

\medskip

We are in the position to prove the main result of this paper, Theorem \ref{theorem:main}.

\begin{proof}[Proof of Theorem \ref{theorem:main}]
First, notice that the space $\N$ of all non-negative integers, endowed with the discrete topology, is homeomorphic to both $\Omega$ and $\Sigma$, so the \v{C}ech--Stone compactifications $\beta\N$, $\beta\Omega$, and $\beta\Sigma$ are mutually homeomorphic. Consequently, by Proposition \ref{prop:fsjnp_product_omega_sigma}, $\beta\N\times\beta\N$ admits a weak* convergent to $0$ sequence $\seqn{\nu_n}$ of finitely supported normalized signed measures with pairwise disjoint supports, contained completely in $\N\times\N$ (as measures $\mu_n$'s defined above on $\beta\Omega\times\beta\Sigma$ have pairwise disjoint supports contained in $\Omega\times\Sigma$).

Let $D$ and $E$ be discrete countable subsets of $K$ and $L$, respectively. Let $\varphi\colon\N\to D$ and $\psi\colon\N\to E$ be bijections. By the Stone Extension Property of $\beta\N$, there are continuous maps $\Phi\colon\beta\N\to K$ and $\Psi\colon\beta\N\to L$ such that $\Phi\rstr\N=\varphi$ and $\Psi\rstr\N=\psi$. For each $n\io$ define a measure $\rho_n$ on $K\times L$ as follows:
\[\rho_n=\sum_{(x,y)\in\supp(\nu_n)}\nu_n\big(\{(x,y)\}\big)\cdot\delta_{(\varphi(x),\psi(y))},\]
it follows that $\big\|\rho_n\big\|=1$ and $\supp\big(\rho_n\big)$ is finite. Since $\seqn{\nu_n}$ weak* converges to $0$, for every $f\in C(K\times L)$ we have:
\[\lim_{n\to\infty}\int_{K\times L}fd\rho_n=\lim_{n\to\infty}\int_{\beta\N\times\beta\N}f(\Phi,\Psi)d\nu_n=0,\]
where $f(\Phi,\Psi)(x,y)=f(\Phi(x),\Psi(y))\in C(\beta\N\times\beta\N)$, so $\seqn{\rho_n}$ is also weak* convergent to $0$.
\end{proof}

%Theorem \ref{theorem:products_fsjnp} yields a strengthening of the result of Khurana \cite[Theorem 2]{Khu78}.

%\begin{corollary}\label{cor:product_ell1_grothendieck}
%For every two infinite compact spaces $K$ and $L$, their product $K\times L$ does not have the $\ell_1$-Grothendieck property. \noproof
%\end{corollary}

%Cembranos \cite[Corollary 2]{Cem84} proved that for a compact space $X$ the space $C(X)$ is Grothendieck if and only if it does not contain any complemented copy of $c_0$ (see also Schachermayer \cite[Proposition 5.3]{Sch82} for a weaker result). This implies Cembranos' result \cite[Corollary 3]{Cem84} stating that the space $C(K\times L)$ admits a complemented copy of $c_0$ for any two infinite compact spaces $K$ and $L$. Corollary \ref{cor:product_ell1_grothendieck}, together with the characterization of 

Theorem \ref{theorem:main} and the aforementioned characterization from \cite{BKS19} of those spaces $C_p(X)$ which contain a complemented copy of the space $(c_0)_p$ imply immediately Corollary \ref{cor:main}. It appears that using this corollary and the Closed Graph Theorem one can easily get Cembranos' and Freniche's Theorem \ref{theorem:cemfre}, as shown in the next proposition.

\begin{proposition}\label{prop:c0p_c0}
Let $X$ be a compact space such that $C_p(X)$ contains a
complemented closed linear subspace $E$ isomorphic to $(c_0)_p$. Then, $E$ with the norm topology of $C(X)$ is complemented in $C(X)$ and isomorphic to
the Banach space $c_0$.
\end{proposition}
\begin{proof}
Let $F$ be a closed linear subspace of $C_p(X)$ such that $C_p(X)=E\oplus
F$. Then,  since the norm topology of $C(X)$ is finer than the
product topology of $C_p(X)$, the spaces $(E,\|\cdot\|)$ and
$(F,\|\cdot\|)$ (i.e. endowed with the inherited norm topology of
$C(X)$) are still closed in $C(X)$ and hence $C(X)=E\oplus F$. It is
enough now to show that $(E,\|\cdot\|)$ is isomorphic to the Banach
space $c_0$. 

Since $(E,\tau_p)$ (i.e. with the inherited product
topology of $C_p(X)$) is isomorphic to $(c_0)_p$, there is a
topology $\tau$ on $E$ stronger than $\tau_p$ and such that $(E,\tau)$ is isomorphic to $c_0$. The identity operator $T\colon(E,\|\cdot\|)\to(E,\tau)$ has closed graph, so it is continuous, and hence $\tau$ is a Banach
space topology on $E$ smaller than the norm topology of $E$. On the
other hand, the identity operator $S\colon(E,\tau)\to(E,\|\cdot\|)$
has closed graph, too, so it is also continuous, and hence the
topology $\tau$ on $E$ is greater than the norm topology of $E$. It
follows that both topologies are equal, and hence
$(E,\|\cdot\|)$ is isomorphic to the Banach space $c_0$.
\end{proof}

%\begin{corollary}[Cembranos--Freniche]\label{cor:product_c0}
%For every two infinite compact spaces $K$ and $L$, the Banach space $C(K\times L)$ contains a complemented copy of $c_0$.
%\end{corollary}

%\begin{proof}
%The result follows from Corollary \ref{cor:main} and Lemma \ref{lemma:c0p_c0}.
%\end{proof}

%Let us note here that the results presented in this paper were also studied and generalized by K\k{a}kol, Marciszewski, Sobota and Zdomskyy \cite{KMSZ20} in the class of pseudocompact spaces.

\begin{thebibliography}{10}

%\bibitem{AK16} F. Albiac, N.J. Kalton, {\em Topics in Banach space theory}, 2nd edition, Graduate Texts in Mathematics 233, Springer, 2016.

\bibitem{BKS19} T. Banakh, J. K\k{a}kol, W. \'Sliwa, {\em Josefson-Nissenzweig property for $C_{p}$-spaces}, RACSAM 113 (2019), 3015--3030.

\bibitem{Bol01} B. Bollob\'as, {\em Random graphs}, 2nd ed., Cambridge Studies in Advanced Mathematics 73, Cambridge University Press, 2001.

\bibitem{Can21} L. Candido, {\em Complementations in $C(K,X)$ and $\ell_\infty(X)$}, preprint (2021), \texttt{arxiv:2104.07152}.

\bibitem{Cem84} P. Cembranos, {\em $C(K, E)$ contains a complemented copy of $c_0$}, Proc. Amer. Math. Soc. 91 (1984), 556--558.

\bibitem{Fre84} F.J. Freniche, {\em Barrelledness of the space of vector valued and simple functions}, Math. Ann. 267 (1984), 479--486.

%\bibitem{Jos75} B. Josefson, {\em Weak sequential convergence in the dual of a Banach space does not imply norm convergence}, Ark. Mat. 13 (1975), 79--89.

%\bibitem{KMSZ20} J. K\k{a}kol, W. Marciszewski, D. Sobota, L. Zdomskyy, {\em On complemented copies of the space $c_0$ in spaces $C_p(X\times Y)$}, to appear in Israel J. Math. (2022), \texttt{arxiv:2007.14723}.

\bibitem{KM20} J. K\k{a}kol, A. Molt\'o, {\em Witnessing the lack of the Grothendieck property in $C(K)$-spaces via convergent sequences}, RACSAM 114:179 (2020).

\bibitem{Khu78} S.S. Khurana, {\em Grothendieck spaces}, Illinois J. Math. 22 (1978), no. 1, 79--80.

\bibitem{Lin67} J. Lindenstrauss, {\em On complemented subspaces of $m$}, Israel J. Math. 5 (1967), 153--156.

%\bibitem{Nis75} A. Nissenzweig, {\em w* sequential convergence}, Israel J. Math. 22 (1975), 266--272.

\bibitem{Sch82} W. Schachermayer, {\em On some classical measure-theoretic theorems for non-sigma-complete Boolean algebras}, Rozpr. Mat. (Diss. Math.) 214 (1982), pp. 34.

\end{thebibliography}

\end{document}
